\section{Konklusion og anbefalinger}
Projektets formål var at udivkle et multiplayer-spil i Unreal Engine med udgangspunkt i de tyske razzia'er af Viborg Katedralskole, med henblik på fejring af bygningens historie samt dets 100 års-jubilæum.

Dette blev opnået ved agil udvikling og brug af en modificeret, scrum-lignende arbejdsprocess. Dette betyder at under hele udviklingsforløbet er der arbejdet iterativt hvor krav, løsninger og design af gameplay elementer løbende er blevet evalueret og tilpasset. 

Brugen af disse metoder bidrog til at skabe struktur i projektet og gav gruppen værktøjer til at håndtere teknisk- og designmæssige udfordringer samt uddeling af arbejdsopgaver samt ugentlige tjek af tilstanden af dem.

Projektet har en tæt knytning til virkeligheden i form af brug af photogrammetry til at give en realistisk repræsentation af skolen, da fokuset især skulle være på selve bygningen. Desuden er grundlaget for spillet tæt knyttet til historiske begivenheder som både er bekræftet gennem egen undersøgelse men også ved samtale med fagfolk hos kunden.

Analysefasen resulterede i en række af user stories, use cases og UML-diagrammer, som sammen skabte et grundlag for systemet videre design og implementation. Disse tog en mere abstrakt ide om spilsystemer og gameplay og strukturede og gjorde de mange forskellige systemer mere overskuelige, samt viste sammenhængen mellem systemerne. 

Designfasen byggede videre fra dette grundlag ved at samle og uddybbe alt i et klassediagram. Desuden takles der her koncepter og krav som analysefasen i mindre detaljegrad har gået over, blandt andet login systemet. 

Implementeringsfasen skete ved brug af OOP koncepter, Blueprints og C++ til udvikling af systemets funktionalitet. Blueprints tillod hurtigt at udvikle gameplay- og UI systemer mens C++ blev anvendt til mere advancerede og performancekritiske funktioner. Til håndtering af multiplayer samt synkronisering af spillerne blev Unreal Engines replication-system brugt i en klient-server-arkitektur. 

Projektet har tydeliggjort flere udfordringer ved udviklingen af multiplayer-spil i Unreal Engine, men også problematikken ved at håndtere store photogrammetry-assets som viste at være mere omfattende end forventet. Dette resulterede i planlagte funktioner, såsom spiller login, spiller data og databaseintegrationen bag ikke blev fuldt implementeret i det endelige produkt. 

Projektet endte ud med et spil som bruger Viborg Katedralskole som fundament for en digital oplevelse med multiplayer funktionalitet. 

Det er diskuterbart hvorvidt projektet fremviser bygningen i et historisk perspektiv, grundet at udgangspunktet af det visuelle aspekt af spillet tager rod i det nuværende udseende og indretning af bygningen. 

Projektet endte ikke med tydeligt at tage udgangspunkt i en bestemt historisk begivenhed, men snarere blot en generel idé om hvordan det kunne have set ud. En mere nøjagtigt repræsentation kræver nok en dybbere indsigt i razzia'en samt inklusioner af assets som mere nøjagtigt afspejler omgivelserne fra tiden. 

Ligeledes endte kundekontakten ikke med elev-samarbejde i forbindelse med detalje- og kreative tilføjelser, som er noget der muligvis kunne undgås ved forsøg på mere tydelig kommunikation med kunden og bedre kravspecifikation ift. hvad eleverne ville skulle arbejde på. 


